\section{AHSME 1984}
        \begin{enumerate}
        	\item (\textit{AHSME 1984}) $\frac{1000^2}{252^2-248^2}$ equals


 $\mathrm{(A) \  }62500 \qquad \mathrm{(B) \  }1000 \qquad \mathrm{(C) \  } 500\qquad \mathrm{(D) \  }250 \qquad \mathrm{(E) \  } \frac{1}{2}$


 
 \textbf{Answer: }C



\textbf{Solution}

We can use \href{https://artofproblemsolving.com/wiki/index.php?title=Difference\_of\_squares}{difference of squares} to factor the \href{https://artofproblemsolving.com/wiki/index.php?title=Denominator}{denominator}, yielding:

 $\frac{1000^2}{252^2-248^2}=\frac{1000^2}{(252-248)(252+248)}=\frac{1000^2}{(4)(500)}=\frac{1000^2}{2000}$.

 We see that the $1000$ in the denominator cancels with one of the $1000$s in the \href{https://artofproblemsolving.com/wiki/index.php?title=Numerator}{numerator}, yielding $500, \boxed{\text{C}}$.

 


	\item (\textit{AHSME 1984}) If $x, y$, and $y-\frac{1}{x}$ are not $0$, then


 
 \[\frac{x-\frac{1}{y}}{y-\frac{1}{x}}\] 
equals


 $\mathrm{(A) \ }1 \qquad \mathrm{(B) \ }\frac{x}{y} \qquad \mathrm{(C) \ } \frac{y}{x}\qquad \mathrm{(D) \ }\frac{x}{y}-\frac{y}{x} \qquad \mathrm{(E) \ } xy-\frac{1}{xy}$


 
 \textbf{Answer: }B



\textbf{Solution}

Multiply the expression by $\frac{xy}{xy}$ to get rid of the fractional \href{https://artofproblemsolving.com/wiki/index.php?title=Numerator}{numerator} and \href{https://artofproblemsolving.com/wiki/index.php?title=Denominator}{denominator}: $\frac{x^2y-x}{xy^2-y}$. This can be \href{https://artofproblemsolving.com/wiki/index.php?title=Factoring}{factored} as $\frac{x(xy-1)}{y(xy-1)}$. The $xy-1$ terms cancel out, leaving $\frac{x}{y}, \boxed{\text{B}}$.

 


	\item (\textit{AHSME 1984}) Let $n$ be the smallest nonprime \href{https://artofproblemsolving.com/wiki/index.php?title=Integer}{integer} greater than $1$ with no \href{https://artofproblemsolving.com/wiki/index.php?title=Prime\_factorization}{prime factor} less than $10$. Then 


 $\mathrm{(A) \ }100<n\leq110 \qquad \mathrm{(B) \ }110<n\leq120 \qquad \mathrm{(C) \ } 120<n\leq130 \qquad \mathrm{(D) \ }130<n\leq140 \qquad \mathrm{(E) \ } 140<n\leq150$


 
 \textbf{Answer: }C



\textbf{Solution}

Since the number isn't \href{https://artofproblemsolving.com/wiki/index.php?title=Prime}{prime}, it is a product of two primes. If the least integer were a product of more than two primes, then one prime could be removed without making the number prime or introducing any prime factors less than $10$. These prime factors must be greater than $10$, so the least prime factor is $11$. Therefore, the least integer is $11^2=121$, which is in $\boxed{\text{C}}$.

 


	\item (\textit{AHSME 1984}) A rectangle intersects a circle as shown: $AB = 4$, $BC = 5$ and $DE = 3$. Then $EF$ equals

                    \begin{figure}[H]
                        centering
                        \includegraphics[width=0.5\textwidth]{Problem/AHSME_I_1984_4_Problem_Figure_1.jpg}
                    \end{figure}
                    


 $\mathrm{(A) \ }6 \qquad \mathrm{(B) \ }7 \qquad \mathrm{(C) \ }\frac{20}{3} \qquad \mathrm{(D) \ }8 \qquad \mathrm{(E) \ }9$


 
 \textbf{Answer: }B



\textbf{Solution}


                    \begin{figure}[H]
                        centering
                        \includegraphics[width=0.5\textwidth]{Solution/AHSME_I_1984_4_Solution_Figure_1.png}
                    \end{figure}
                    

 Draw $BE$ and $CF$, forming a \href{https://artofproblemsolving.com/wiki/index.php?title=Trapezoid}{trapezoid}. Since it's cyclic, this trapezoid must be \href{https://artofproblemsolving.com/wiki/index.php?title=Isosceles\_trapezoid}{isosceles}. Also, drop \href{https://artofproblemsolving.com/wiki/index.php?title=Altitude}{altitudes} from $E$ to $AC$, $B$ to $DF$, and $C$ to $DF$, and let the feet of these altitudes be $G$, $H$, and $I$ respectively. $AGED$ is a \href{https://artofproblemsolving.com/wiki/index.php?title=Rectangle}{rectangle} since it has $4$ \href{https://artofproblemsolving.com/wiki/index.php?title=Right\_angle}{right angles}. Therefore, $AG=DE=3$, and $GB=4-3=1$. By the same logic, $GBHE$ is also a rectangle, and $EH=GB=1$. $BH=CI$ since they're both altitudes to a trapezoid, and $BE=CF$ since the trapezoid is isosceles. Therefore, $\triangle BHE \cong \triangle CIF$ by HL \href{https://artofproblemsolving.com/wiki/index.php?title=Congruence}{congruence}, so $IF=EH=1$. Also, $BCIH$ is a rectangle from $4$ right angles, and $HI=BC=5$. Therefore, $EF=EH+HI+IF=1+5+1=\boxed{7}$.

 


	\item (\textit{AHSME 1984}) The largest \href{https://artofproblemsolving.com/wiki/index.php?title=Integer}{integer} $n$ for which $n^{200}<5^{300}$ is


 $\mathrm{(A) \ }8 \qquad \mathrm{(B) \ }9 \qquad \mathrm{(C) \ } 10 \qquad \mathrm{(D) \ }11 \qquad \mathrm{(E) \ } 12$


 
 \textbf{Answer: }D



\textbf{Solution}

Since both sides are positive, we can take the $100th$ root of both sides to find the largest integer $n$ such that $n^2<5^3$. Fortunately, this is simple to evaluate: $5^3=125$, and the largest \href{https://artofproblemsolving.com/wiki/index.php?title=Perfect\_square}{square} less than $125$ is $11^2=121$, so the largest $n$ is $11, \boxed{\text{D}}$.

 


	\item (\textit{AHSME 1984}) In a certain school, there are three times as many boys as girls and nine times as many girls as teachers. Using the letters $b, g, t$ to represent the number of boys, girls and teachers, respectively, then the total number of boys, girls and teachers can be represented by the \href{https://artofproblemsolving.com/wiki/index.php?title=Expression}{expression}


 $\mathrm{(A) \ }31b \qquad \mathrm{(B) \ }\frac{37}{27}b \qquad \mathrm{(C) \ } 13g \qquad \mathrm{(D) \ }\frac{37}{27}g \qquad \mathrm{(E) \ } \frac{37}{27}t$


 
 \textbf{Answer: }B



\textbf{Solution}

From the given, we have $3g=b$ and $9t=g$, or $t=\frac{g}{9}$. The sum of these, in terms of $g$, is $3g+g+\frac{g}{9}$, or, with a common \href{https://artofproblemsolving.com/wiki/index.php?title=Denominator}{denominator}, $\frac{37g}{9}$. We can see that this isn't one of the choices. So we write it in terms of $b$. We can see from the first \href{https://artofproblemsolving.com/wiki/index.php?title=Equation}{equation} that $g=\frac{b}{3}$, so substituting this into the expression yields $\frac{37b}{27}, \boxed{\text{B}}$.

 


	\item (\textit{AHSME 1984}) When Dave walks to school, he averages $90$ steps per minute, each of his steps $75$ cm long. It takes him $16$ minutes to get to school. His brother, Jack, going to the same school by the same route, averages $100$ steps per minute, but his steps are only $60$ cm long. How long does it take Jack to get to school?


 $\mathrm{(A) \ }14 \frac{2}{9} \text{ minutes} \qquad \mathrm{(B) \ }15 \text{ minutes}\qquad \mathrm{(C) \ } 18 \text{ minutes}\qquad \mathrm{(D) \ }20 \text{ minutes} \qquad \mathrm{(E) \ } 22 \frac{2}{9} \text{ minutes}$


 
 \textbf{Answer: }C



\textbf{Solution}

At $90$ steps per minute and $75 cm$ per step, Dave walks at a rate of $90\times75$ cm per minute, and with $16$ minutes, the distance Dave walks to school is $90\times75\times16$. Also, at $100$ steps per minute and $60$ cm per step, Jack walks at a rate of $100\times60$ cm per minute. Jack must walk $90\times75\times16$ cm, so it takes him $\frac{90\times75\times16}{100\times60}$ minutes. Canceling some of the factors, this comes out to $18 \text{minutes}, \boxed{\text{C}}$.

 


	\item (\textit{AHSME 1984}) Figure $ABCD$ is a \href{https://artofproblemsolving.com/wiki/index.php?title=Trapezoid}{trapezoid} with $AB||DC$, $AB=5$, $BC=3\sqrt{2}$, $\angle BCD=45^\circ$ and $\angle CDA=60^\circ$. The length of $DC$ is


 
                    \begin{figure}[H]
                        centering
                        \includegraphics[width=0.5\textwidth]{Problem/AHSME_I_1984_8_Problem_Figure_1.jpg}
                    \end{figure}
                    


 $\mathrm{(A) \ }7+\frac{2}{3}\sqrt{3} \qquad \mathrm{(B) \ }8 \qquad \mathrm{(C) \ } 9 \frac{1}{2} \qquad \mathrm{(D) \ }8+\sqrt{3} \qquad \mathrm{(E) \ } 8+3\sqrt{3}$


 
 \textbf{Answer: }D



\textbf{Solution}


                    \begin{figure}[H]
                        centering
                        \includegraphics[width=0.5\textwidth]{Solution/AHSME_I_1984_8_Solution_Figure_1.png}
                    \end{figure}
                    Drop \href{https://artofproblemsolving.com/wiki/index.php?title=Perpendicular}{perpendiculars} from $A$ and $B$ to $CD$, thus forming a \href{https://artofproblemsolving.com/wiki/index.php?title=Right\_triangle\#Special\_right\_triangles}{30-60-90 right triangle}, a \href{https://artofproblemsolving.com/wiki/index.php?title=Right\_triangle\#Special\_right\_triangles}{45-45-90  right triangle}, and a \href{https://artofproblemsolving.com/wiki/index.php?title=Rectangle}{rectangle}, as shown in the diagram. Since $ABEF$ is a rectangle, we have $EF=AB=5$. Also, from the properties of 45-45-90 right triangles, we have $EB=EC=\frac{BC}{\sqrt{2}}=\frac{3\sqrt{2}}{\sqrt{2}}=3$. Again since $ABEF$ is a rectangle, $AF=3$, too. From the properties of 30-60-90 right triangles, we have $DF=\frac{AF}{\sqrt{3}}=\frac{3}{\sqrt{3}}=\sqrt{3}$. Therefore, $CD=CE+EF+FD=3+5+\sqrt{3}=8+\sqrt{3}, \boxed{\text{D}}$.

 


	\item (\textit{AHSME 1984}) The number of \href{https://artofproblemsolving.com/wiki/index.php?title=Digit}{digits} in $4^{16}5^{25}$ (when written in the usual \href{https://artofproblemsolving.com/wiki/index.php?title=Base\_number}{base} $10$ form) is


 $\mathrm{(A) \ }31 \qquad \mathrm{(B) \ }30 \qquad \mathrm{(C) \ } 29 \qquad \mathrm{(D) \ }28 \qquad \mathrm{(E) \ } 27$


 
 \textbf{Answer: }D



\textbf{Solution}

We can rewrite this as $2^{32}5^{25}$. We can also combine some of the factors to introduce factors of $10$, whose digit count is simple to evaluate because it simply adds $0$s. Thus, we have $2^{32}5^{25}=2^72^{25}5^{25}=2^710^{25}$. We can see that this final number is $2^7$ with $25$ $0$s annexed onto it. $2^7=128$, which has $3$ digits, so the entire number has $25+3=28$ digits, $\boxed{\text{D}}$.

 


	\item (\textit{AHSME 1984}) Four \href{https://artofproblemsolving.com/wiki/index.php?title=Complex\_numbers}{complex numbers} lie at the \href{https://artofproblemsolving.com/wiki/index.php?title=Vertices}{vertices} of a \href{https://artofproblemsolving.com/wiki/index.php?title=Square}{square} in the \href{https://artofproblemsolving.com/wiki/index.php?title=Complex\_plane}{complex plane}. Three of the numbers are $1+2i, {-2}+i$ and ${-1}-2i$. The fourth number is 


 $\mathrm{(A) \ }2+i \qquad \mathrm{(B) \ }2-i \qquad \mathrm{(C) \ } 1-2i \qquad \mathrm{(D) \ }{-1}+2i \qquad \mathrm{(E) \ } {-2}-i$


 
 \textbf{Answer: }B



\textbf{Solution}

Perhaps the easiest way to attack this is to transfer this to the \href{https://artofproblemsolving.com/wiki/index.php?title=Cartesian\_plane}{Cartesian plane}. The points then would be $(1, 2), (-2, 1),$ and $(-1, -2)$, assuming the real axis was horizontal. Let these points be $A, B$ and $C$, respectively. The remaining point is then the intersection of the following \href{https://artofproblemsolving.com/wiki/index.php?title=Perpendicular}{perpendicular} \href{https://artofproblemsolving.com/wiki/index.php?title=Line}{lines}:The one perpenicular to $AB$ and passing through $A$ and the one perpendicular to $BC$ and passing through $C$.The \href{https://artofproblemsolving.com/wiki/index.php?title=Slope}{slope} of the first line is the negative \href{https://artofproblemsolving.com/wiki/index.php?title=Reciprocal}{reciprocal} of the slope of the line through $AB$, which, using the slope formula, is $\frac{1}{3}$, so the slope of the perpendicular line is $-3$. It passes through $(1,2)$, so the equation of the line in point slope form is $y-2=-3(x-1)$, or $y=-3x+5$. Similarly, the slope of the second line is $\frac{-1}{-3}=\frac{1}{3}$, and, since it passes through $(-1,-2)$, its equation is $y+2=\frac{1}{3}(x+1)$, or $y=\frac{1}{3}x-\frac{5}{3}$. To find the intersection, we have $-3x+5=y=\frac{1}{3}x-\frac{5}{3}$, and solving for $x$ yields $x=2$. Plugging this back into the equation yields $y=-1$, so the remaining point in the Cartesian plane is $(2, -1)$, and in the complex plane is $2-i, \boxed{\text{B}}$.

 


 NICE! You can also just graph and use pythag.Or get the solution by sheer obviousness (I'm serious)

 


	\item (\textit{AHSME 1984}) A calculator has a key that replaces the displayed entry with its \href{https://artofproblemsolving.com/wiki/index.php?title=Perfect\_square}{square}, and another key which replaces the displayed entry with its \href{https://artofproblemsolving.com/wiki/index.php?title=Reciprocal}{reciprocal}. Let $y$ be the final result when one starts with an entry $x\not=0$ and alternately squares and reciprocates $n$ times each. Assuming the calculator is completely accurate (e.g. no roundoff or overflow), then $y$ equals


 $\mathrm{(A) \ }x^{((-2)^n)} \qquad \mathrm{(B) \ }x^{2n} \qquad \mathrm{(C) \ } x^{-2n} \qquad \mathrm{(D) \ }x^{-(2^n)} \qquad \mathrm{(E) \ } x^{((-1)^n2n)}$


 
 \textbf{Answer: }A



\textbf{Solution}

Notice that taking the reciprocal of a number is equivalent to raising the number to the $-1st$ power. Therefore, taking the reciprocal of a number $x$ and then squaring it is $(x^{-1})^2=x^{-2}$. Also, since multiplication is commutative, this is equivalent to squaring the number then taking the reciprocal of it. Doing this $n$ times equals $\left(\left(x^{-2}\right)^{-2}\right)^{\cdots})) (\text{n times})=x^{(-2)^n}, \boxed{\text{A}}$.

 


	\item (\textit{AHSME 1984}) If the \href{https://artofproblemsolving.com/wiki/index.php?title=Sequence}{sequence} $\{a_n\}$ is defined by


 
 \[a_1=2,\]

 \[a_{n+1}=a_n+2n \qquad (n \geq 1),\]


 then $a_{100}$ equals


 $\mathrm{(A) \ }9900 \qquad \mathrm{(B) \ }9902 \qquad \mathrm{(C) \ } 9904 \qquad \mathrm{(D) \ }10100 \qquad \mathrm{(E) \ } 10102$


 
 \textbf{Answer: }B



\textbf{Solution}

We begin to evaluate the first couple of terms of the sequence, hoping to find a pattern:$2, 4, 8, 14, 22, ....$. We notice that the difference between succesive terms of the sequence are $2, 4, 6, 8, ....$, a clear pattern. We can see that this pattern continues infinitely because of the recursive definition: each term is the previous term plus the next even number. Therefore, since the differences of consecutive terms form an arithmetic sequence, then the terms satisfy a \href{https://artofproblemsolving.com/wiki/index.php?title=Quadratic}{quadratic}, specifically, the one that contains the points $(1, 2), (2, 4),$, and $(3, 8)$. Let the quadratic be $f(x)=ax^2+bx+c$, so:

 $a+b+c=2$ \textbf{(1)}

 $4a+2b+c=4$ \textbf{(2)}

 $9a+3b+c=8$ \textbf{(3)}

 Subtracting \textbf{(1)} from \textbf{(2)} and \textbf{(2)} from \textbf{(3)} yields the two-variable \href{https://artofproblemsolving.com/wiki/index.php?title=System\_of\_equations}{system of equations}

 $3a+b=2$ \textbf{(4)}

 $5a+b=4$ \textbf{(5)}

 We can subtract \textbf{(4)} from \textbf{(5)} to find that $2a=2$, so $a=1$. Substituting this back in yields $b=-1$, and substituting these back into one of the original equations yields $c=2$, so the closed form for the terms is

 $f(x)=x^2-x+2$, or

 $a_n=n^2-n+2$.

 Substituting in $n=100$ yields $a_{100}=100^2-100+2=9902, \boxed{\text{B}}$.

 \textbf{Alternate Solution} Term $a_{100}$ can be viewed as the first term, $2$, plus the arithmetic series $2 + 4 +\cdots + 198$.

 By summing the numbers that form the arithmetic sequence with $a_{1}$,  we get  $2 + \dfrac{2 + 198}{2}\cdot99 = 9902$.

 


	\item (\textit{AHSME 1984}) $\frac{2\sqrt{6}}{\sqrt{2}+\sqrt{3}+\sqrt{5}}$ equals


 $\mathrm{(A) \ }\sqrt{2}+\sqrt{3}-\sqrt{5} \qquad \mathrm{(B) \ }4-\sqrt{2}-\sqrt{3} \qquad \mathrm{(C) \ } \sqrt{2}+\sqrt{3}+\sqrt{6}-5 \qquad$


 $\mathrm{(D) \ }\frac{1}{2}(\sqrt{2}+\sqrt{5}-\sqrt{3}) \qquad \mathrm{(E) \ } \frac{1}{3}(\sqrt{3}+\sqrt{5}-\sqrt{2})$


 
 \textbf{Answer: }A



\textbf{Solution}

Multiply the \href{https://artofproblemsolving.com/wiki/index.php?title=Numerator}{numerator} and the \href{https://artofproblemsolving.com/wiki/index.php?title=Denominator}{denominator} by $\sqrt{2}-(\sqrt{3}+\sqrt{5})$ to get $\frac{2\sqrt{6}(\sqrt{2}-\sqrt{3}-\sqrt{5})}{(\sqrt{2}+(\sqrt{3}+\sqrt{5}))(\sqrt{2}-(\sqrt{3}+\sqrt{5})}$ 

 $=\frac{2\sqrt{12}-2\sqrt{18}-2\sqrt{30}}{2-(\sqrt{3}+\sqrt{5})^2}$ 

 $=\frac{4\sqrt{3}-6\sqrt{2}-2\sqrt{30}}{-6-2\sqrt{15}}$ 

 $=\frac{2\sqrt{3}-3\sqrt{2}-\sqrt{30}}{-3-\sqrt{15}}$.

 Now to get rid of the $\sqrt{15}$ in the denominator. Multiply the numerator and denominator by $-3+\sqrt{15}$ to get

 $\frac{(2\sqrt{3}-3\sqrt{2}-\sqrt{30})(-3+\sqrt{15})}{(-3-\sqrt{15})(-3+\sqrt{15})}$

 
$=\frac{-6\sqrt{3}+9\sqrt{2}+3\sqrt{30}+2\sqrt{45}-3\sqrt{30}-\sqrt{450}}{9-15}$

 
$=\frac{-6\sqrt{3}+9\sqrt{2}+3\sqrt{30}+6\sqrt{5}-3\sqrt{30}-15\sqrt{2}}{-6}$

 
$=\sqrt{2}+\sqrt{3}-\sqrt{5}, \boxed{\text{A}}$.

 



\textbf{Solution 2}

Multiply the numerator and denominator by $\sqrt{2}+\sqrt{3}-\sqrt{5}$. We get $\frac{(2\sqrt{6})(\sqrt{2}+\sqrt{3}-\sqrt{5})}{(\sqrt{2}+\sqrt{3}+\sqrt{5})(\sqrt{2}+\sqrt{3}-\sqrt{5})}$

 
 \[\frac{(2\sqrt{6})(\sqrt{2}+\sqrt{3}-\sqrt{5})}{2+3-5+2\sqrt{6}+\sqrt{10}-\sqrt{10}+\sqrt{15}-\sqrt{15}}=\frac{(2\sqrt{6})(\sqrt{2}+\sqrt{3}-\sqrt{5})}{2\sqrt{6}}=\sqrt{2}+\sqrt{3}-\sqrt{5}\]$\boxed{A}$

 


	\item (\textit{AHSME 1984}) The product of all real \href{https://artofproblemsolving.com/wiki/index.php?title=Root\_(polynomial)}{roots} of the equation $x^{\log_{10}{x}}=10$ is


 $\mathrm{(A) \ }1 \qquad \mathrm{(B) \ }-1 \qquad \mathrm{(C) \ } 10 \qquad \mathrm{(D) \ }10^{-1} \qquad \mathrm{(E) \ } \text{None of these}$


 
 \textbf{Answer: }A



\textbf{Solution}

Take the \href{https://artofproblemsolving.com/wiki/index.php?title=Logarithm}{logarithm} base $x$ of both sides to get $\log_{10}x=\log_x10$. Use the \href{https://artofproblemsolving.com/wiki/index.php?title=Change\_of\_base\_formula}{change of base formula} to get $\log_{10}x=\frac{\log_{10}10}{\log_{10}x}=\frac{1}{\log_{10}x}$. Therefore, $(\log_{10}x)^2=1$ and $\log_{10}x=\pm1$. From $\log_{10}x=1$, we have $x=10^1=10$ and from $\log_{10}x=-1$ we have $x=10^{-1}=\frac{1}{10}$. The product is thus $10\times\frac{1}{10}=1, \boxed{\text{A}}$.

 


	\item (\textit{AHSME 1984}) If $\sin{2x}\sin{3x}=\cos{2x}\cos{3x}$, then one value for $x$ is


 $\mathrm{(A) \ }18^\circ \qquad \mathrm{(B) \ }30^\circ \qquad \mathrm{(C) \ } 36^\circ \qquad \mathrm{(D) \ }45^\circ \qquad \mathrm{(E) \ } 60^\circ$


 
 \textbf{Answer: }A



\textbf{Solution}

We divide both sides of the \href{https://artofproblemsolving.com/wiki/index.php?title=Equation}{equation} by $\cos{2x}\times\cos{3x}$ to get $\frac{\sin{2x}\times\sin{3x}}{\cos{2x}\times\cos{3x}}=1$, or $\tan{2x}\times\tan{3x}=1$.

 This looks a lot like the formula relating the slopes of two \href{https://artofproblemsolving.com/wiki/index.php?title=Perpendicular}{perpendicular} \href{https://artofproblemsolving.com/wiki/index.php?title=Line}{lines}, which is $m_1\times m_2=-1$, where $m_1$ and $m_2$ are the \href{https://artofproblemsolving.com/wiki/index.php?title=Slope}{slopes}. It's made even more relatable by the fact that the \href{https://artofproblemsolving.com/wiki/index.php?title=Tangent}{tangent} of an \href{https://artofproblemsolving.com/wiki/index.php?title=Angle}{angle} can be defined by the slope of the line that makes that angle with the \href{https://artofproblemsolving.com/wiki/index.php?title=X-axis}{x-axis}.

 We can make this look even more like the slope formula by multiplying both sides by $-1$:

 $\tan{2x}\times-\tan{3x}=-1$, and using the \href{https://artofproblemsolving.com/wiki/index.php?title=Trigonometric\_identity}{trigonometric identity} $-\tan{x}=\tan{-x}$, we have $\tan{2x}\times\tan{-3x}=-1$.

 Now it's time for a diagram:

 
                    \begin{figure}[H]
                        centering
                        \includegraphics[width=0.5\textwidth]{Solution/AHSME_I_1984_15_Solution_Figure_1.png}
                    \end{figure}
                    Since the product of the two slopes, $\tan{2x}$ and $\tan{-3x}$, is $-1$, the lines are perpendicular, and the \href{https://artofproblemsolving.com/wiki/index.php?title=Angle}{angle} between them is $\frac{\pi}{2}$. The angle between them is also $2x+3x=5x$, so $5x=\frac{\pi}{2}$ and $x=\frac{\pi}{10}$, or $18^\circ, \boxed{\text{A}}$.

 NOTE: To show that $18^\circ$ is not the only solution, we can also set $5x$ equal to another angle measure congruent to $\frac{\pi}{2}$, such as $\frac{5\pi}{2}$, yielding another solution as $\frac{\pi}{2}=90^\circ$, which clearly is a solution to the equation.

 



\textbf{Solution 2}

We can simply try the answers. We quickly see that $18^\circ$ works, since $\sin(36^\circ)\sin(54^\circ)=\cos(54^\circ)\cos(36^\circ)=\cos(36^\circ)\cos(54^\circ)$ by the identity $\cos(x)=\sin(90^\circ-x)$

 



\textbf{Solution 3}

Start by subtracting $\sin{2x}\sin{3x}$ from both sides to get $\cos{2x}\cos{3x}-\sin{2x}\sin{3x}=0$. We recognize that this is of the form $\cos{a}\cos{b}-\sin{a}\sin{b}=\cos{(a+b)}$, so $\cos{2x}\cos{3x}-\sin{2x}\sin{3x}=\cos{(2x+3x)}=0$. $\cos{90^\circ}=0$, so $x=\boxed{18^\circ}$.

 


	\item (\textit{AHSME 1984}) The \href{https://artofproblemsolving.com/wiki/index.php?title=Function}{function} $f(x)$ satisfies $f(2+x)=f(2-x)$ for all real numbers $x$. If the equation $f(x)=0$ has exactly four distinct real \href{https://artofproblemsolving.com/wiki/index.php?title=Root\_(polynomial)}{roots}, then the sum of these roots is


 $\mathrm{(A) \ }0 \qquad \mathrm{(B) \ }2 \qquad \mathrm{(C) \ } 4 \qquad \mathrm{(D) \ }6 \qquad \mathrm{(E) \ } 8$


 
 \textbf{Answer: }E



\textbf{Solution}

Let one of the roots be $r_1$. Also, define $x$ such that $2+x=r_1$. Thus, we have $f(2+x)=f(r_1)=0$ and $f(2+x)=f(2-x)$. Therefore, we have $f(2-x)=0$, and $2-x$ is also a root. Let this root be $r_2$. The sum $r_1+r_2=2+x+2-x=4$. Similarly, we can let $r_3$ be a root and define $y$ such that $2+y=r_3$, and we will find $2-y$ is also a root, say, $r_4$, so $r_3+r_4=2+y+2-y=4$. Therefore, $r_1+r_2+r_3+r_4=4+4=8, \boxed{\text{E}}$.

 



\textbf{Solution 2}

The graph of this function is symmetric around 2. Therefore, two roots will be greater than $2$ and the other two roots will be less than $2$. These four roots are symmetric around $2$, so the average of the four roots is $2$. Then, the sum is $2\times4=8$.

 


	\item (\textit{AHSME 1984}) A \href{https://artofproblemsolving.com/wiki/index.php?title=Right\_triangle}{right triangle} $ABC$ with \href{https://artofproblemsolving.com/wiki/index.php?title=Hypotenuse}{hypotenuse} $AB$ has side $AC=15$. \href{https://artofproblemsolving.com/wiki/index.php?title=Altitude}{Altitude} $CH$ divides $AB$ into segments $AH$ and $HB$, with $HB=16$. The \href{https://artofproblemsolving.com/wiki/index.php?title=Area}{area} of $\triangle ABC$ is


 
                    \begin{figure}[H]
                        centering
                        \includegraphics[width=0.5\textwidth]{Problem/AHSME_I_1984_17_Problem_Figure_1.jpg}
                    \end{figure}
                    


 $\mathrm{(A) \ }120 \qquad \mathrm{(B) \ }144 \qquad \mathrm{(C) \ } 150 \qquad \mathrm{(D) \ }216 \qquad \mathrm{(E) \ } 144\sqrt{5}$


 
 \textbf{Answer: }C



\textbf{Solution}


                    \begin{figure}[H]
                        centering
                        \includegraphics[width=0.5\textwidth]{Solution/AHSME_I_1984_17_Solution_Figure_1.png}
                    \end{figure}
                    $AHC\sim ACB$ by $AA$, so $\frac{AH}{AC}=\frac{AC}{AH+16}$. Since $AC=15$, we have $\frac{AH}{15}=\frac{15}{AH+16}$. Cross mutliplying, we have $AH(AH+16)=225$. Solving this \href{https://artofproblemsolving.com/wiki/index.php?title=Quadratic}{quadratic} yields $AH=9$. Also, $AHC\sim CHB$, so $\frac{AH} {HC}=\frac{HC}{HB}$. Substituting in known values, we have $\frac{9}{HC}=\frac{HC}{16}$, so $HC^2=144$ and $HC=12$. 

 The area of $\triangle ABC$ is $\frac{1}{2}(AB)(HC)=\frac{1}{2}(25)(12)=150, \boxed{\text{C}}$.

 


	\item (\textit{AHSME 1984}) A point $(x, y)$ is to be chosen in the \href{https://artofproblemsolving.com/wiki/index.php?title=Coordinate\_plane}{coordinate plane} so that it is equally distant from the $x$-axis, the $y$-axis, and the \href{https://artofproblemsolving.com/wiki/index.php?title=Line}{line} $x+y=2$. Then $x$ is


 $\mathrm{(A) \ }\sqrt{2}-1 \qquad \mathrm{(B) \ }\frac{1}{2} \qquad \mathrm{(C) \ } 2-\sqrt{2} \qquad \mathrm{(D) \ }1 \qquad \mathrm{(E) \ } \text{not uniquely determined}$


 
 \textbf{Answer: }E



\textbf{Solution}

The $x$-axis and $y$-axis intersect at $(0,0)$, while the line $x+y=2$ intersects the $x$-axis at $(2,0)$ and the $y$-axis at $(0,2)$.  Let $A (0,2)$, $B (2,0)$, and $C (0,0)$ be these three points.  Then the incenter and the three excenters of $\triangle ABC$ must all be equidistant from all three of these lines.  Since the $B$-excenter has a negative $x$-coordinate while the incenter and other two excenters have positive $x$-coordinates, $x$ is $\boxed{(\mathbf{E})\ \mathrm{not\ uniquely\ determined}}$.

 -j314andrews

 


	\item (\textit{AHSME 1984}) A box contains $11$ balls, numbered $1, 2, 3,\ldots,11$. If $6$ balls are drawn simultaneously at random, what is the \href{https://artofproblemsolving.com/wiki/index.php?title=Probability}{probability} that the sum of the numbers on the balls drawn is odd?


 $\mathrm{(A) \ }\frac{100}{231} \qquad \mathrm{(B) \ }\frac{115}{231} \qquad \mathrm{(C) \ } \frac{1}{2} \qquad \mathrm{(D) \ }\frac{118}{231} \qquad \mathrm{(E) \ } \frac{6}{11}$


 
 \textbf{Answer: }D



\textbf{Solution}

There are exactly $3$ ways the sum can be odd: exactly $5$ are odd, exactly $3$ are odd, or exactly $1$ is odd. There are exactly $6$ odd numbers in the set. For the first case, there are ${6\choose5}\times{5\choose1}=30$ possibilities. For the second case, there are ${6\choose3}\times{5\choose3}=200$ possibilities. For the third case, there are ${6\choose1}\times{5\choose5}=6$ possibilities. In each case, we're choosing out of the $6$ odd numbers which ones will be picked, and multiplying that by out of the $5$ even numbers, which one of those will be picked so that there are $6$ numbers. In total,there are $236$ desirable possibilities out of a total of ${11\choose6}=462$ possibilities. Thus, the answer is $\frac{236}{462}=\frac{118}{231}, \boxed{\text{D}}$.

 


	\item (\textit{AHSME 1984}) The number of the distinct solutions of the \href{https://artofproblemsolving.com/wiki/index.php?title=Equation}{equation} $|{x-|2x+1|}|=3$ is


 $\mathrm{(A) \ }0 \qquad \mathrm{(B) \ }1 \qquad \mathrm{(C) \ } 2 \qquad \mathrm{(D) \ }3 \qquad \mathrm{(E) \ } 4$


 
 \textbf{Answer: }C



\textbf{Solution}

We can create a tree of possibilities, progressively eliminating the absolute value signs by creating different cases.

 
                    \begin{figure}[H]
                        centering
                        \includegraphics[width=0.5\textwidth]{Solution/AHSME_I_1984_20_Solution_Figure_1.png}
                    \end{figure}
                    

 So we have $4$ possible solutions: $-2, \frac{2}{3}, 2,$ and $-\frac{4}{3}$. Checking for extraneous solutions, we find that the only ones that work are $2$ and $-\frac{4}{3}$, so there are $2$ solutions, $\boxed{\text{C}}$.

 


	\item (\textit{AHSME 1984}) The number of triples $(a, b, c)$ of \href{https://artofproblemsolving.com/wiki/index.php?title=Natural\_numbers}{positive integers} which satisfy the simultaneous \href{https://artofproblemsolving.com/wiki/index.php?title=Equation}{equations}


 
 \[ab+bc=44,\]

 \[ac+bc=23,\]


 is


 $\mathrm{(A) \ }0 \qquad \mathrm{(B) \ }1 \qquad \mathrm{(C) \ } 2 \qquad \mathrm{(D) \ }3 \qquad \mathrm{(E) \ } 4$


 
 \textbf{Answer: }C



\textbf{Solution}

We can \href{https://artofproblemsolving.com/wiki/index.php?title=Factoring}{factor} the second equation to get $c(a+b)=23$, so we see that $c$ must be a \href{https://artofproblemsolving.com/wiki/index.php?title=Factor}{factor} of $23$, and since this is \href{https://artofproblemsolving.com/wiki/index.php?title=Prime}{prime}, $c=1$ or $c=23$. However, if $c=23$, then $a+b=1$, which is impossible for the field of positive integers. Therefore, $c=1$ for all possible solutions. Substituting this into the original equations gives

 $ab+b=44$

 and 

 $a+b=23$.

 From the second equation, $a=23-b$, and substituting this into the first equation yields $b(23-b)+b=44$, or $b^2-24b+44=0$. Factoring this gives $(b-2)(b-22)=0$, so $b=2$ or $b=22$. Both of these yield integer solutions for $a$, giving $a=21$ or $a=1$, respectively.Therefore, the only solutions are $(21, 2, 1)$ and $(1, 22, 1)$, yielding $2$ solutions, $\boxed{\text{C}}$.

 


	\item (\textit{AHSME 1984}) Let $a$ and $c$ be fixed \href{https://artofproblemsolving.com/wiki/index.php?title=Natural\_numbers}{positive numbers}. For each \href{https://artofproblemsolving.com/wiki/index.php?title=Real\_number}{real number} $t$ let $(x_t, y_t)$ be the \href{https://artofproblemsolving.com/wiki/index.php?title=Vertex}{vertex} of the \href{https://artofproblemsolving.com/wiki/index.php?title=Parabola}{parabola} $y=ax^2+tx+c$. If the set of the vertices $(x_t, y_t)$ for all real values of $t$ is graphed on the \href{https://artofproblemsolving.com/wiki/index.php?title=Cartesian\_plane}{plane}, the \href{https://artofproblemsolving.com/wiki/index.php?title=Graph}{graph} is


 $\mathrm{(A) \ } \text{a straight line} \qquad \mathrm{(B) \ } \text{a parabola} \qquad \mathrm{(C) \ } \text{part, but not all, of a parabola} \qquad \mathrm{(D) \ } \text{one branch of a hyperbola} \qquad \mathrm{(E) \ } \text{none of these}$


 
 \textbf{Answer: }B



\textbf{Solution}

The $x$-coordinate of the vertex of the parabola $ax^2 + tx + c$ is $-\frac{t}{2a}$, so $x_t=-\frac{t}{2a}$. Substituting yields $y_t=-\frac{t^2}{4a^2}+c$. Notice that $y_t=-\frac{t^2}{4a^2}+c=-a(-\frac{t}{2a})^2+c=-ax_t^2+c$, so all of the vertices are on a parabola. However, we have only showed that all of the points in the locus of vertices are on a parabola, we have not shown whether or not all points on the parabola are on the locus. Assume we are given an $x_t$ on the parabola. $-\frac{t}{2a}=x_t$, $t = -2ax_t$, so a unique $t$, and therefore a unique vertex, is determined for each point on the parabola. We therefore conclude that every point in the locus is on the parabola and every point on the parabola is in the locus, and the graph of the locus is the same as the graph of the parabola, $\boxed{\text{B}}$.

 


	\item (\textit{AHSME 1984}) $\frac{\sin{10^\circ}+\sin{20^\circ}}{\cos{10^\circ}+\cos{20^\circ}}$ equals


 $\mathrm{(A) \ }\tan{10^\circ}+\tan{20^\circ} \qquad \mathrm{(B) \ }\tan{30^\circ} \qquad \mathrm{(C) \ } \frac{1}{2}(\tan{10^\circ}+\tan{20^\circ}) \qquad \mathrm{(D) \ }\tan{15^\circ} \qquad \mathrm{(E) \ } \frac{1}{4}\tan{60^\circ}$


 
 \textbf{Answer: }D



\textbf{Solution}

From the \href{https://artofproblemsolving.com/wiki/index.php?title=Trigonometric\_identities\#Prosthaphaeresis\_Identities}{sum-to-product formulas}, we have $\sin{a}+\sin{b}=2\sin{\frac{a+b}{2}}\cos{\frac{a-b}{2}}$. So, $\sin{10^\circ}+\sin{20^\circ}=2\sin{15^\circ}\cos{5^\circ}$. Also from the sum-to-product formulas, we have $\cos{a}+\cos{b}=2\cos{\frac{a+b}{2}}\cos{\frac{a-b}{2}}$, so $\cos{10^\circ}+\cos{20^\circ}=2\cos{15^\circ}\cos{5^\circ}$. Substituting these back into the original fraction, we have $\frac{\sin{10^\circ}+\sin{20^\circ}}{\cos{10^\circ}+\cos{20^\circ}}=\frac{2\sin{15^\circ}\cos{5^\circ}}{2\cos{15^\circ}\cos{5^\circ}}=\frac{\sin{15^\circ}}{\cos{15^\circ}}=\tan{15^\circ}, \boxed{\text{D}}$.

 


	\item (\textit{AHSME 1984}) If $a$ and $b$ are positive \href{https://artofproblemsolving.com/wiki/index.php?title=Real\_numbers}{real numbers} and each of the \href{https://artofproblemsolving.com/wiki/index.php?title=Equation}{equations} $x^2+ax+2b=0$ and $x^2+2bx+a=0$ has real \href{https://artofproblemsolving.com/wiki/index.php?title=Root\_(polynomial)}{roots}, then the smallest possible value of $a+b$ is


 $\mathrm{(A) \ }2 \qquad \mathrm{(B) \ }3 \qquad \mathrm{(C) \ } 4 \qquad \mathrm{(D) \ }5 \qquad \mathrm{(E) \ } 6$


 
 \textbf{Answer: }E



\textbf{Solution}

Since both of the equations have real roots, both of their \href{https://artofproblemsolving.com/wiki/index.php?title=Discriminant}{discriminants} are nonnegative. Therefore, we have 

 $a^2-4(1)(2b)=a^2-8b\geq0\implies a^2\geq8b$ from the first equation, and 

 $(2b)^2-4(1)(a)=4b^2-4a\geq0\implies 4b^2\geq4a\implies b^2\geq a$ from the second.

 We can square the second equation to get $b^4\geq a^2$, and combining this with the first one gives $b^4\geq a^2\geq8b$, so $b^4\geq8b$. We can divide both sides by $b$, since it is positive, and take the cubed root of that to get $b\geq2$. Therefore, we have $a^2\geq8b\geq8(2)=16$, and since $a$ is positive, we can take the \href{https://artofproblemsolving.com/wiki/index.php?title=Square\_root}{square root} of this to get $a\geq4$. Therefore, $a+b\geq 2+4=6$, and the smallest possible value is $6, \boxed{\text{E}}$. 

 


	\item (\textit{AHSME 1984}) The total \href{https://artofproblemsolving.com/wiki/index.php?title=Area}{area} of all the \href{https://artofproblemsolving.com/wiki/index.php?title=Faces}{faces} of a \href{https://artofproblemsolving.com/wiki/index.php?title=Rectangular\_prism}{rectangular solid} is $22 \ \text{cm}^2$, and the total length of all its \href{https://artofproblemsolving.com/wiki/index.php?title=Edges}{edges} is $24 \ \text{cm}$. Then the length in cm of any one of its \href{https://artofproblemsolving.com/wiki/index.php?title=Diagonal\#Polyhedra}{interior diagonals} is


 $\mathrm{(A) \ }\sqrt{11} \qquad \mathrm{(B) \ }\sqrt{12} \qquad \mathrm{(C) \ } \sqrt{13}\qquad \mathrm{(D) \ }\sqrt{14} \qquad \mathrm{(E) \ } \text{not uniquely determined}$


 
 \textbf{Answer: }D



\textbf{Solution}

Let the edge lengths be $a, b$, and $c$. Therefore, the total area of all its faces is $2ab+2bc+2ac$. Therefore, $2ab+2bc+2ac=22$ and $ab+bc+ac=11$. Also, the total lengths of all of its edges is $4a+4b+4c$, so $4a+4b+4c=24$, and $a+b+c=6$. Therefore, we have:

 $ab+bc+ac=11$

 and

 $a+b+c=6$.

 The length of the interior diagonal is $\sqrt{a^2+b^2+c^2}$, so if we can find $a^2+b^2+c^2$, we can find the diagonal. We square the second equation above to introduce an $a^2+b^2+c^2$:

 $(a+b+c)^2=a^2+b^2+c^2+2ab+2bc+2ac=36$.

 However, we already know that $2ab+2bc+2ac=22$, so $a^2+b^2+c^2+22=36\implies a^2+b^2+c^2=14\implies \sqrt{a^2+b^2+c^2}=\sqrt{14}, \boxed{\text{D}}$.

 


	\item (\textit{AHSME 1984}) In the \href{https://artofproblemsolving.com/wiki/index.php?title=Obtuse\_triangle}{obtuse triangle} $ABC$, $AM=MB$, $MD\perp BC$, $EC\perp BC$. If the \href{https://artofproblemsolving.com/wiki/index.php?title=Area}{area} of $\triangle ABC$ is $24$, then the area of $\triangle BED$ is


 
                    \begin{figure}[H]
                        centering
                        \includegraphics[width=0.5\textwidth]{Problem/AHSME_I_1984_26_Problem_Figure_1.jpg}
                    \end{figure}
                    


 $\mathrm{(A) \ }9 \qquad \mathrm{(B) \ }12 \qquad \mathrm{(C) \ } 15 \qquad \mathrm{(D) \ }18 \qquad \mathrm{(E) \ } \text{not uniquely determined}$


 
 \textbf{Answer: }B



\textbf{Solution 1}


                    \begin{figure}[H]
                        centering
                        \includegraphics[width=0.5\textwidth]{Solution/AHSME_I_1984_26_Solution_Figure_1.png}
                    \end{figure}
                    

 We let side $BC$ have length $a$, $AB$ have length $c$, and $\angle ABC$ have angle measure $\beta$. We then have that

 
 \[[ABC]=24=\frac{AB\cdot BC\sin{\angle ABC}}{2}=\frac{ac\sin{\beta}}{2}\]

 Now I shall find the lengths of $BD$ and $CE$ in terms of the defined variables. Note that $M$ is defined to be the midpoint of $AB$, so $BM=\frac{c}{2}$. We can then use trigonometric manipulation on triangle $BDM$ to get that $BD=\frac{c\cos{\beta}}{2}$. We can also use trig manipulation on $BCE$ to get that $CE=a\tan{\beta}$.

 Now note that $CE$ is the height of triangle $BDE$ originating from vertex $E$, so we have that

 $[BED]=\frac{BD\cdot CE}{2}=\frac{ac\cos{\beta}\tan{\beta}}{4}=\frac{ac\sin{\beta}}{4}$

 However, this is simply half the area of triangle $ABC$, so $[BED]=12$, which makes $\boxed{\textbf{B}}$ the correct answer.

 



\textbf{Solution 2}


                    \begin{figure}[H]
                        centering
                        \includegraphics[width=0.5\textwidth]{Solution/AHSME_I_1984_26_Solution_Figure_2.png}
                    \end{figure}
                    

 Connect points $M$ and $C$ to make $\overline{MC}$. $[BED]$ can be expressed as the sum of two smaller triangles:$[BED] = [MDB]+[EMD]$. Since $\triangle EMD$ and $\triangle MDC$ share the same base $MD$ and height along the same line, we have $[EMD] = [MDC]$.Therefore,$[BED] = [MDB] + [MDC] = [MBC] = \frac{1}{2}*[ABC] = \frac{1}{2}*24 = 12$. The answer is $\boxed{C}$.

 -BiegeTomcat72

 


	\item (\textit{AHSME 1984}) In $\triangle ABC$, $D$ is on $AC$ and $F$ is on $BC$. Also, $AB\perp AC$, $AF\perp BC$, and $BD=DC=FC=1$. Find $AC$.


 
                    \begin{figure}[H]
                        centering
                        \includegraphics[width=0.5\textwidth]{Problem/AHSME_I_1984_27_Problem_Figure_1.jpg}
                    \end{figure}
                    


 $\mathrm{(A) \ }\sqrt{2} \qquad \mathrm{(B) \ }\sqrt{3} \qquad \mathrm{(C) \ } \sqrt[3]{2} \qquad \mathrm{(D) \ }\sqrt[3]{3} \qquad \mathrm{(E) \ } \sqrt[4]{3}$


 
 \textbf{Answer: }C



\textbf{Solution}


                    \begin{figure}[H]
                        centering
                        \includegraphics[width=0.5\textwidth]{Solution/AHSME_I_1984_27_Solution_Figure_1.png}
                    \end{figure}
                    

 Let $AC=x$ and $BF=y$. We have $AFC\sim BFA$ by AA, so $\frac{AF}{FC}=\frac{BF}{AF}$. Substituting in known values gives $\frac{AF}{1}=\frac{y}{AF}$, so $AF=\sqrt{y}$. Also, $AD=x-1$, and using the \href{https://artofproblemsolving.com/wiki/index.php?title=Pythagorean\_Theorem}{Pythagorean Theorem} on $\triangle ABD$, we have $AB^2+(x-1)^2=1^2$, so $AB=\sqrt{2x-x^2}$. Using the Pythagorean Theorem on $\triangle AFC$ gives $y+1=x^2$, or $y=x^2-1$. Now, we use the Pythagorean Theorem on $\triangle AFB$ to get $y^2+y=2x-x^2$. Substituting $y=x^2-1$ into this gives $(x^2-1)^2+x^2-1=2x-x^2$, or $x^4-2x^2+1+x^2-1=2x-x^2$. Simplifying this and moving all of the terms to one side gives $x^4-2x=0$, and since $x\not=0$, we can divide by $x$ to get $x^3-2=0$, from which we find that $x=\sqrt[3]{2}, \boxed{\text{C}}$.

 


	\item (\textit{AHSME 1984}) The number of distinct pairs of \href{https://artofproblemsolving.com/wiki/index.php?title=Integers}{integers} $(x, y)$ such that $0<x<y$ and $\sqrt{1984}=\sqrt{x}+\sqrt{y}$ is


 $\mathrm{(A) \ }0 \qquad \mathrm{(B) \ }1 \qquad \mathrm{(C) \ } 3 \qquad \mathrm{(D) \ }4 \qquad \mathrm{(E) \ } 7$


 
 \textbf{Answer: }C



\textbf{Solution}

We can simplify $\sqrt{1984}$ to $8\sqrt{31}$. Therefore, the only solutions are $a\sqrt{31}+b\sqrt{31}$ such that $a+b=8$ and $0<a<b$. The only solutions to this are $a=1, b=7; a=2, b=6; a=3, b=5$. Each of these gives distinct pairs of $(x, y)$, so there are $3$ pairs, $\boxed{\text{C}}$.

 


	\item (\textit{AHSME 1984}) Find the largest value for $\frac{y}{x}$ for pairs of \href{https://artofproblemsolving.com/wiki/index.php?title=Real\_numbers}{real numbers} $(x, y)$ which satisfy $(x-3)^2+(y-3)^2=6$.


 $\mathrm{(A) \ }3+2\sqrt{2} \qquad \mathrm{(B) \ }2+\sqrt{3} \qquad \mathrm{(C) \ } 3\sqrt{3} \qquad \mathrm{(D) \ }6 \qquad \mathrm{(E) \ } 6+2\sqrt{3}$


 
 \textbf{Answer: }A



\textbf{Solution}

Let $\frac{y}{x}=k$, so that $y=kx$. Substituting this into the given \href{https://artofproblemsolving.com/wiki/index.php?title=Equation}{equation} $(x-3)^2+(y-3)^2=6$ yields $(x-3)^2+(kx-3)^2=6$. Multiplying this out and forming it into a \href{https://artofproblemsolving.com/wiki/index.php?title=Quadratic}{quadratic} yields $(1+k^2)x^2+(-6-6k)x+12=0$. 

 We want $x$ to be a \href{https://artofproblemsolving.com/wiki/index.php?title=Real\_number}{real number}, so we must have the \href{https://artofproblemsolving.com/wiki/index.php?title=Discriminant}{discriminant} $\geq0$. The discriminant is $(-6-6k)^2-4(1+k^2)(12)=36k^2+72k+36-48-48k^2=-12k^2+72k-12$. Therefore, we must have $-12k^2+72k-12\geq0$, or $k^2-6k+1\leq0$. The \href{https://artofproblemsolving.com/wiki/index.php?title=Root\_(polynomial)}{roots} of this quadratic, using the \href{https://artofproblemsolving.com/wiki/index.php?title=Quadratic\_formula}{quadratic formula}, are $3\pm2\sqrt{2}$, so the quadratic can be \href{https://artofproblemsolving.com/wiki/index.php?title=Factoring}{factored} as $(k-(3-2\sqrt{2}))(k-(3+2\sqrt{2}))\leq0$. We can now separate this into $3$ cases:

 \textbf{Case 1: $k<3-2\sqrt{2}$}Then, both \href{https://artofproblemsolving.com/wiki/index.php?title=Term&action=edit&redlink=1}{terms} in the factored quadratic are negative, so the inequality doesn't hold.

 \textbf{Case 2: $3-2\sqrt{2}<k<3+2\sqrt{2}$}Then, the first term is positive and the second is negative and the second is positive, so the inequality holds.

 \textbf{Case 3: $k>3+2\sqrt{2}$}Then, both terms are positive, so the inequality doesn't hold.

 Also, when $k=3-2\sqrt{2}$ or $k=3+2\sqrt{2}$, the equality holds.

 Therefore, we must have $3-2\sqrt{2}\leq k\leq3+2\sqrt{2}$, and the maximum value of $k=\frac{y}{x}$ is ${3+2\sqrt{2}}, \boxed{\text{A}}$.

 



\textbf{Solution 2}

The equation represents a circle of radius $\sqrt{6}$ centered at $A=(3,3)$. To find the maximal $k$ with $y=kx$ is equivalent to finding the maximum slope of a line passing through the origin $O$ and intersecting the circle. The steepest such line is tangent to the circle at some point $X$. We have $XA=\sqrt{6}$, $OA=\sqrt{18}$, $\angle OXA = 90$ because the line is tangent to the circle. Using the pythagorean theorem, we have $OX=\sqrt{12}$. 

 The slope we are looking for is equivalent to $\tan (\theta + 45)$ where $\angle AOX = \theta$. Using tangent addition,

 
 \[\tan (\theta + 45)= \frac{\tan \theta + \tan 45}{1-\tan\theta\tan 45} = \frac{\frac{1}{\sqrt{2}}+1}{1-\frac{1}{\sqrt{2}}}=3+2\sqrt{2}\]

 So $\boxed{A}$ is the answer

 



\textbf{Solution 3}

Following the steps of Solution 2, we get $OX=\sqrt{12}$. Let $X = (x, y)$. Based on the distance to the origin and to the centre of the circle

 
 \[x^2 + y^2 = 12\]

 
 \[(x - 3)^2 + (y - 3)^2 = 6\]

 Expanding the latter and substituting the former into the equation gives $30 - 6(x + y) = 6$. $x + y = 4$. We can square this and remove the 1st equation to get $2xy = 4$, or $xy = 2$.

 Substitute $y = 4 - x$ to get

 
 \[x(4 - x) = 2\]

 
 \[x^2 - 4x + 2 = 0\]

 
 \[x = 2 \pm \sqrt{2}\]

 Since $x$ and $y$ are symmetric, we want $x$ to be the lower number to maximize $k$, so $x = 2 - \sqrt{2}$ and $y = 2 + \sqrt{2}$. 

 
 \[\frac{y}{x} = \frac{2 + \sqrt{2}}{2 - \sqrt{2}}\]

 
 \[\frac{y}{x} = \frac{(2 + \sqrt{2})^2}{2}\]

 
 \[\frac{y}{x} = 3 + 2\sqrt{2}\]

 


	\item (\textit{AHSME 1984}) For any \href{https://artofproblemsolving.com/wiki/index.php?title=Complex\_number}{complex number} $w=a+bi$, $|w|$ is defined to be the \href{https://artofproblemsolving.com/wiki/index.php?title=Real\_number}{real number} $\sqrt{a^2+b^2}$. If $w=\cos40^\circ+i\sin40^\circ$, then
$|w+2w^2+3w^3+...+9w^9|^{-1}$ equals


 $\mathrm{(A) \ }\frac{1}{9}\sin40^\circ \qquad \mathrm{(B) \ }\frac{2}{9}\sin20^\circ \qquad \mathrm{(C) \ } \frac{1}{9}\cos40^\circ \qquad \mathrm{(D) \ }\frac{1}{18}\cos20^\circ \qquad \mathrm{(E) \ } \text{none of these}$


 
 \textbf{Answer: }B



\textbf{Solution 1}

Let $S=w+2w^2+3w^3+...+9w^9$. Note that

 
 \[S=\sum_{i=1}^{9}\sum_{j=i}^{9} w^j\]

 Now we multiply $S$ by $1-w$:

 
 \[S(1-w)=\sum_{i=1}^{9}\sum_{j=i}^{9} w^j(1-w)\]

 By the geometric series formula, $\sum_{j=i}^{9} w^j(1-w)$ is simply $w^i-w^{10}$. Therefore
 \[S(1-w)=\sum_{i=1}^{9} w^i-w^{10}=(w+w^2+\cdots +w^9)-9w^{10}\]

 A simple application of \href{https://artofproblemsolving.com/wiki/index.php?title=De\_Moivre\%27s\_Theorem}{De Moivre's Theorem} shows that $w$ is a ninth root of unity ($w^9=1$), so

 
 \[(w+w^2+\cdots +w^9)(1-w)=w-w^{10}=w(1-w^{9})=0\]

 This shows that $S=\frac{-9w^{10}}{1-w}$. Note that $w^{10}=w\cdot w^9=w$, so $S=\frac{-9w}{1-w}$. 

 It's not hard to show that $\left|S\right|^{-1}=\left|S^{-1}\right|=\left|-S^{-1}\right|$, so the number we seek is equal to $\left|\frac{1-w}{9w}\right|$.

 Now we plug $w$ into the fraction:

 
 \[\frac{1-w}{9w}=\frac{(1-\cos{40^{\circ}})-i\sin{40^{\circ}}}{9\cos{40}+9i\sin{40^{\circ}}}\]

 We multiply the numerator and denominator by $9\cos{40^{\circ}}-9i\sin{40^{\circ}}$ and simplify to get

 
 \[\frac{1-w}{9w}=\frac{(\cos{40^{\circ}}-i\sin{40^{\circ}})((1-\cos{40^{\circ}})-i\sin{40^{\circ}})}{9}\]

 
 \[=\frac{\cos{40^{\circ}}-\cos^2{40^{\circ}}-\sin^2{40^{\circ}}+i(-\sin{40^{\circ}}+\sin{40^{\circ}}\cos{40^{\circ}}-\sin{40^{\circ}}\cos{40^{\circ}})}{9}\]

 
 \[=\frac{(\cos{40^{\circ}}-1)-i\sin{40^{\circ}}}{9}\]

 The absolute value of this is

 
 \[\left|\frac{1-w}{9w}\right|=\frac{1}{9}\sqrt{(\cos{40^{\circ}}-1)^2+\sin^2{40^{\circ}}}=\frac{1}{9}\sqrt{1-2\cos{40^{\circ}}+\cos^2{40^{\circ}}+\sin^2{40^{\circ}}}=\frac{\sqrt{2}}{9}\sqrt{1-\cos{40^{\circ}}}\]

 Note that, from double angle formulas, $\cos{40^{\circ}}=\cos^2{20^{\circ}}-\sin^2{20^{\circ}}$, so $1-\cos{40^{\circ}}=\cos^2{20^{\circ}}+\sin^2{20^{\circ}}-(\cos^2{20^{\circ}}-\sin^2{20^{\circ}})=2\sin^2{20^{\circ}}$. Therefore

 
 \[\left|\frac{1-w}{9w}\right|=\frac{\sqrt{2}}{9} \sqrt{2\sin^2{20^{\circ}}}\]

 
 \[=\frac{2}{9}\sin{20^{\circ}}\]

 Therefore the correct answer is $\boxed{\textbf{(B)}}$.

 



\textbf{Solution 2}

Notice that $w = \cos40^\circ+i\sin40^\circ = \cos\left(\frac{2\pi}{9}\right)+i\sin\left(\frac{2\pi}{9}\right)$ is the primitive root of the equation $x^9 = 1$.

 Consider the simple geometric progression,

 
 \[f(x) = 1 + x + x^2 +\cdots + x^9 = \frac{x^{10} - 1}{x - 1}\]

 Differentiating this equation (using the \href{https://artofproblemsolving.com/wiki/index.php/Quotient\_Rule}{Quotient Rule}), we obtain the following equation,

 
 \[f'(x) = 1 + 2x + 3x^2 + \cdots + 9x^8 = \frac{9x^{10} - 10x^9 + 1}{(x - 1)^2}\]

 Now, we can evaluate the given expression,

 
 \[w + 2w^2 + 3w^3 + \cdots + 9w^9 = w \cdot (1+2w+3w^2+...+9w^8)\]
 \[= w\cdot f'(w)\]
 \[= w\cdot\frac{9w^{10} - 10w^9 + 1}{(w - 1)^2}\]
 \[= \frac{9w}{w-1} \; (\because w^9 = 1)\]

 Finally, we can calculate the given modulus,

 
 \[\lvert w+2w^2+3w^3+...+9w^9 \rvert^{-1} = \lvert \frac{1}{9}\cdot\left(1 - w^8) \right) \rvert\]
 \[= \lvert \frac{1}{9}\cdot\left(1 - \cos\left(320^\circ\right)+i\sin\left(320^\circ\right)) \right) \rvert\]
 \[= \frac{1}{9}\sqrt{\left(1 - \cos\left(320^\circ\right)\right)^2+\left(\sin\left(320^\circ\right)\right)^2}\]
 \[= \frac{1}{9}\sqrt{2 - 2\cos\left(320^\circ\right)}\]
 \[= \frac{2}{9}\sin\left(160^\circ\right)\]
 \[= \boxed{\frac{2}{9}\sin\left(20^\circ\right)}\]

 


\end{enumerate}